Let $\xb = (x_1,x_2,\ldots)$. For $n \ge 2$ and $D \subseteq [n-1]$, the \emph{fundamental quasisymmetric function} $Q_{n,D}(\xb)$ is given by
\[
Q_{n,D}(\xb) = \sum_{\substack{a_1 \le a_2 \le \cdots \le a_n, \\ a_i = a_{i+1} \Rightarrow i \not\in D}} x_{a_1} x_{a_2} \cdots x_{a_n}.
\]

The \emph{Schur function} $s_\lambda(\xb)$ indexed by $\lambda$ is written in terms of quasisymmetric functions as follows:
\[
s_\lambda(\xb) = \sum_{T \in \SYT(\lambda)} Q_{n,\Des(T)}(\xb).
\]
Schur functions form a basis for the ring of \emph{symmetric functions} $\Lambda(\xb) \subseteq \ZZ[[x_1,x_2,\ldots]]$, the subspace consisting of all power series of bounded degree which are unchanged by permuting variable indices. Another important basis for the ring of symmetric functions is formed by the \emph{power sum} symmetric functions $p_\lambda(\xb)$, which are given by
\[
p_d(\xb) = x_1^d + x_2^d + \cdots, \qquad p_\lambda(\xb) = p_{\lambda_1}(\xb)p_{\lambda_2}(\xb) \cdots
\]
for any $d \ge 1$ and any $\lambda = (\lambda_1 \ge \lambda_2 \ge \cdots > 0)$.

\subsection{Representation theory of \texorpdfstring{$\GL(\CC^N)$}{GL(C\^{}N)}}
We next provide some background about the representation theory of the general linear group $\GL(\CC^N)$ of invertible linear transformations of $\CC^N$. We discuss its Schur--Weyl duality with the symmetric group $S_k$ --- a complete exposition of which can be found, for instance, in \cite{Ful} --- and then recount the relevant facts about its character theory and (bi)graded representations.

Let $N \ge 1$. A \emph{$\GL(\CC^N)$-representation} $E$ is a $\CC$-vector space equipped with an action of $\GL(\CC^N)$. All representations in this paper will be polynomial. The irreducible representations of $\GL(\CC^N)$ are the \emph{Schur modules} $V^\lambda$, and are indexed by integer partitions $\lambda = (\lambda_1 \ge \cdots \ge \lambda_N \ge 0)$ with at most $N$ non-zero parts. 

By the Schur--Weyl duality between the representation theory of $\GL(\CC^N)$ and the symmetric group $S_k$, the irreducible representations $V^\lambda$ may be constructed as follows. The irreducible representations of $S_k$ are the \emph{Specht modules} $S^\lambda$, which are indexed by partitions $\lambda \vdash k$. The \emph{Frobenius characteristic map} $\FrobCh$ is an isomorphism between the space of class functions on $S_k$ and the space of homogeneous degree $k$ symmetric functions, defined by $\FrobCh(S^\lambda) = s_\lambda(\xb)$ for any $\lambda \vdash k$. Letting $V = \CC^N$, the space $V^{\otimes k}$ admits a natural $\GL(V) \times S_k$-action, where $\GL(V)$ acts diagonally on $V^{\otimes k}$ on the left, and $S_k$ acts on the right by permuting tensor factors. For any $S_k$-module $M$, the \emph{Schur--Weyl dual} of $M$ is the $\GL(V)$-module
\[
E(M) = V^{\otimes k} \otimes_{\CC S_k} M.
\]
Then the irreducible representations of $\GL(V)$ are given as the Schur--Weyl duals of all $S^\lambda$ for $\ell(\lambda) \le N$.

Upon fixing a basis for $\CC^N$, we may identify $\GL(\CC^N)$ with the group of invertible $N \times N$ matrices over $\CC$. The \emph{character} of a $\GL(\CC^N)$-representation $E$ is the trace of the action of the diagonal matrix $\mathrm{diag}(x_1,\ldots,x_N)$ on $E$, denoted $\Ch(E;\xb_N)$. The character $\Ch(E;\xb_N)$ is a polynomial in the variables $\xb_N = x_1,\ldots,x_N$, and the character of the Schur module $V^\lambda$ is
\[
\Ch(V^\lambda;\xb_N) = s_\lambda(x_1,\ldots,x_N,0,0,\ldots) \in \ZZ[x_1,\ldots,x_N],
\]
the Schur polynomial in $N$ variables. We then have that
\[
\lim_{N \to \infty} \Ch(V^\lambda;\xb_N) = s_\lambda(\xb) = \FrobCh(S^\lambda),
\]
so in practice we will often let $N \to \infty$ and omit reference to the underlying vector space. By the $\ZZ$-linear independence of Schur functions, computing the irreducible decomposition of a $\GL(V)$- or $S_k$-module is equivalent to computing the expansion of its character (resp. Frobenius characteristic) in the Schur basis.

A \emph{graded} (resp. \emph{bigraded}) $\GL(\CC^N)$-representation is a graded (resp. doubly graded) $\CC$-vector space $E = \bigoplus_n E_n$ (resp. $F = \bigoplus_{n,m} F_{n,m}$) equipped with a $\GL(\CC^N)$-action, such that the graded pieces $E_n$ (resp. $F_{n,m}$) are themselves $\GL(\CC^N)$-representations. The \emph{graded characters} $\grch(E;\xb;q)$ and $\grch(F;\xb,q,t)$ of $E$ and $F$ (resp.) are defined as
\[
\grch(E;\xb_N,q) = \sum_n \Ch(E_n;\xb_N)q^n, \quad \grch(F;\xb_N,q,t) = \sum_{n,m} \Ch(F_{n,m};\xb_N)q^nt^m.
\]

We will later consider representations of $\GL(\CC^N) \oplus \GL(\CC^M)$. After fixing bases for $\CC^N$ and $\CC^M$, we may consider the action of $\mathrm{diag}(x_1,\ldots,x_N,y_1,\ldots,y_M)$ on a $\GL(\CC^N) \oplus \GL(\CC^M)$-representation $E$. As above, the \emph{character} $\Ch(E;\xb_N,\yb_M)$ of $E$ is the trace of this action on $E$, which is a polynomial in $\CC[x_1,\ldots,x_N;y_1,\ldots,y_M]$, symmetric in the $x$- and $y$-variables separately. If $E = \bigoplus_{n,m} E_{n,m}$ is a bigraded $\GL(\CC^N) \oplus \GL(\CC^M)$-representation, then its graded character is defined as
\[
\grch(E;\xb_N,\yb_M,q,t) = \sum_{n,m} \Ch(E_{n,m};\xb_N,\yb_M)q^nt^m.
\]
\subsection{The classical case of Thrall's problem}\label{sec:Thrall-classical}
%Recall that a Lie algebra $L$ over $\CC$ is a $\CC$-vector space equipped with an operation $[-,-]$ which is bilinear, anticommutative, and satisfies the Jacobi identity.
We now define Thrall's problem and recall some of the progress that has been made on it. The tensor algebra decomposition considered by Thrall is described in \eqref{PBW-decomp}, and the higher Lie modules $\Lcal_\lambda$ are defined in \Cref{def:lie-mod}. We then outline the three results related to Thrall's problem which we seek to generalize: Klyachko's identification of the Schur--Weyl dual of $\Lcal_n$ (\Cref{thm:kly}), Brandt's formula (\Cref{thm:brandt}), and Kra{\'s}kiewicz--Weyman's Schur decomposition of $\Ch(\Lcal_n;\xb)$ (\Cref{thm:KW}).

Let $V$ be a $\CC$-vector space with basis $X = \{ x_1,\ldots,x_N \}$, and let $\Trm(V) = \bigoplus_{n \ge 0} V^{\otimes n}$ denote its tensor algebra, which is naturally a graded $\GL(\CC^N)$-module. The \emph{free Lie algebra} $\Lcal(V)$ is the Lie subalgebra of $\Trm(V)$ generated by $X$. Then $\Lcal(V) = \bigoplus_{n \ge 1} \Lcal_n(V)$ inherits a grading from $\Trm(V)$, where $\Lcal_n(V) \coloneqq \Lcal(V) \cap V^{\otimes n}$, endowing $\Lcal(V)$ with the structure of a graded $\GL(\CC^N)$-module. See \cite{Reu} for more background on the free Lie algebra.

It is well-known that the universal enveloping algebra $U(\Lcal(V))$ of $\Lcal(V)$ is naturally isomorphic to the tensor algebra $\Trm(V)$. On the other hand, the Poincar\'e--Birkhoff--Witt theorem gives $U(\Lcal(V)) \cong S(\Lcal(V))$, where $\Srm(W) = \bigoplus_{m \ge 0} \Srm^m(W)$ denotes the symmetric algebra of a complex vector space $W$. Hence we have the following decomposition of $\Trm(V)$ as a $\GL(\CC^N)$-module:
\begin{align}
\Trm(V) &\cong U(\Lcal(V)) \cong \Srm(\Lcal(V)) \cong \bigotimes_{n \ge 1} \Srm(\Lcal_n(V)) = \bigotimes_{n \ge 1} \left( \bigoplus_{m \ge 0} \Srm^m(\Lcal_n(V)) \right) \nonumber \\
&\cong \bigoplus_{\lambda = (1^{m_1}2^{m_2} \cdots)} \Srm^{m_1}(\Lcal_1(V)) \otimes \Srm^{m_2}(\Lcal_2(V)) \otimes \cdots. \label{PBW-decomp}
\end{align}
We more carefully derive an analogous decomposition of the tensor algebra $\Trmtil(V)$ of a super vector space in \Cref{thm:super-Thrall}. The isomorphism in \eqref{PBW-decomp} motivates the following definition.
\begin{defi}\label{def:lie-mod}
For $\lambda = (1^{m_1}2^{m_2} \cdots)$, the \emph{higher Lie module} $\Lcal_\lambda(V)$ is:
\[
\Lcal_\lambda(V) = \Srm^{m_1}(\Lcal_1(V)) \otimes \Srm^{m_2}(\Lcal_2(V)) \otimes \cdots.
\]
\end{defi}
Thus the higher Lie modules $\{ \Lcal_\lambda(V) : \lambda \in \mathrm{Par} \}$ yield a decomposition of $\Trm(V)$ as a $\GL(\CC^N)$-module, and it is therefore natural to ask for their irreducible decompositions. This problem was originally posed by Thrall \cite{Thrall}. By taking characters, an equivalent formulation of Thrall's problem is as follows.
\begin{prob}[Thrall's problem]
  Determine the coefficients $a_\mu \in \ZZ_{\ge 0}$ in the Schur expansion of $\Ch(\Lcal_\lambda(V);\xb_N)$:
  \[
  \Ch(\Lcal_\lambda(V);\xb_N) = \sum_\mu a_\mu s_\mu(\xb_N).
  \]
\end{prob}
As noted above, we may take $N \to \infty$ and work instead with symmetric functions; in this case we will simply write $\Lcal_\lambda$ for the higher Lie module, omitting the reference to the underlying vector space.

By \eqref{PBW-decomp} and the Littlewood--Richardson rule, it would suffice to determine the irreducible components of $\Lcal_{(a^b)} = S^b(\Lcal_a)$ for the purposes of Thrall's problem. On the level of characters, we have
\[
\Ch(\Lcal_{(a^b)};\xb_N) = h_b[\Ch(\Lcal_a)],
\]
where the right-hand side denotes the plethysm of $h_b$ with $\Ch(\Lcal_a)$ (see, for instance, \cite[Appendix~2]{Stan}). Thrall's problem remains open in general, although the single-row case follows from work of Klyachko \cite{Kly} and Kra{\'s}kiewicz--Weyman \cite{KW} by identifying the Schur--Weyl dual of $\Lcal_n$.

Let $\pi_n = (12 \cdots n) \in S_n$, and let $\omega_n$ denote a primitive $n$-th root of unity. Let $C_n = \langle \pi_n \rangle \le S_n$ denote the cyclic subgroup of order $n$ generated by $\pi_n$. The irreducible characters $\chi^1, \ldots, \chi^n$ of $C_n$ are given by $\chi^k(\pi_n) = \omega_n^k$. Klyachko proved that the Schur--Weyl dual of $\Lcal_n$ is obtained by inducing the representation $\chi^1$ of $C_n$ up to $S_n$:
\begin{thm}[\cite{Kly}]\label{thm:kly}
For any $n \ge 1$, the Schur--Weyl dual of $\Lcal_n$ is $\chi^1 \uparrow_{C_n}^{S_n}$; that is,
\[
\Ch(\Lcal_n;\xb) = \FrobCh(\chi^1 \uparrow_{C_n}^{S_n};\xb).
\]
\end{thm}
This result is equivalent to an expansion for $\Ch(\Lcal_n;\xb)$ in the power sum basis, which was also found directly by Brandt:
\begin{thm}[\cite{Brandt}]\label{thm:brandt}
For any $n \ge 1$,
\[
\Ch(\Lcal_n;\xb) = \frac{1}{n} \sum_{d \mid n} \mu(d) p_d(\xb)^{\frac{n}{d}},
\]
where $\mu(-)$ denotes the M{\"o}bius function.
\end{thm}
Kra{\'s}kiewicz--Weyman leverage the above expansion to determine the expansion of $\Ch(\Lcal_n;\xb)$ in the Schur basis.
\begin{thm}[\cite{KW}]\label{thm:KW}
For any $n \ge 1$, we have
\[
\Ch(\Lcal_n;\xb) = \sum_{\mu \vdash n} a_{\mu,1} s_\mu(\xb),
\]
where $a_{\mu,1} = |\{ T \in \SYT(\mu) : \maj(T) \equiv_n 1 \}|$.
\end{thm}
Outside of the single-row case, the only other known case of Thrall's problem is when $\lambda = (2^b)$, in which case $\Ch(\Lcal_{(2^b)})$ is given by a known plethystic identity:
\[
\Ch(\Lcal_{(2^b)}) = h_b[e_2] = \sum_{\mu} s_\mu,
\]
where the sum is over all $\mu \vdash 2b$ whose columns all have even length (see e.g. \cite[Ex.~I.8.6(b)]{Mac}).

\section{Super Thrall's problem}\label{sec:Thrall}
We now wish to extend Thrall's problem to the setting of free Lie superalgebras. In \Cref{sec:lie-super} we recall some facts about (free) Lie superalgebras, and then obtain a super generalization of Brandt's formula, \Cref{thm:super-brandt}. In \Cref{sec:super-Thrall} we identify the analogue of the higher Lie modules for the free Lie superalgebra, and state the super generalization of Thrall's problem in \Cref{thm:super-Thrall}.

\subsection{Free Lie superalgebras}\label{sec:lie-super}
Here we state the necessary background on Lie superalgebras, and then describe the structure of the free Lie superalgebra $\Ltil(V)$ as a bigraded $\GL(\CC^N) \oplus \GL(\CC^M)$-module. A result of Petrogradsky (\Cref{thm:free-super-lie-char}) determines the $\GL(\CC^N) \oplus \GL(\CC^M)$-character of $\Ltil(V)$, and we use this to obtain a character formula for the bigraded components of $\Ltil(V)$ in \Cref{thm:super-bi-brandt}, from which we obtain a proof of \Cref{thm:super-brandt}. We work over $\CC$ throughout.

A \textit{superalgebra} $\Atil$ is a $\ZZ/2$-graded vector space $\Atil = \Atil_0 \oplus \Atil_1$ with an associative, bilinear multiplication with a unit and satisfying $\Atil_i \Atil_j \subset \Atil_{i+j}$, where the indices are taken modulo $2$. For $x \in \Atil_i$, write $|x| = i$. The \textit{super commutator} on $\Atil$ is defined by extending
  \[ [x, y] \coloneqq xy - (-1)^{|x||y|} yx \]
bilinearly. It satisfies 
\begin{itemize}
\item[S1.] $[x,y] = -(-1)^{|x||y|}[y,x]$,
\item[S2.] $(-1)^{|x||z|}[x,[y,z]] + (-1)^{|z||y|}[z,[x,y]] + (-1)^{|y||x|}[y,[z,x]] = 0$.
\end{itemize}

\begin{exam}
    The \textit{tensor superalgebra} of $V = V_0 \oplus V_1 = \CC^N \oplus \CC^M$ is
      \[ \Trmtil(V) \coloneqq \bigoplus_{d=0}^\infty (V_0 \oplus V_1)^{\otimes d} \]
    with the natural concatenation product. It has a bigrading where $\Trmtil_{n, m}(V)$ is spanned by tensors with $n$ factors from $V_0$ and $m$ factors from $V_1$. The $\ZZ/2$-grading is given by
      \[ \Trmtil_0(V) = \bigoplus_{n, m \geq 0} \Trmtil_{n, 2m}(V) \qquad\text{and}\qquad \Trmtil_1(V) = \bigoplus_{n, m \geq 0} \Trmtil_{n, 2m+1}(V). \]
    Both gradings respect the natural $\GL(\CC^N) \oplus \GL(\CC^M)$-action.
\end{exam}

\begin{exam}
    The \textit{symmetric superalgebra} of $V = V_0 \oplus V_1 = \CC^N \oplus \CC^M$ is
      \[ \Stil(V) \coloneqq S(V_0) \otimes \bigwedge(V_1) \]
    where $S(W)$ is the symmetric algebra and $\bigwedge(W)$ is the exterior algebra of the vector space $W$. As with $\Trmtil(V)$, it is a bigraded and $\ZZ/2$-graded $\GL(\CC^N) \oplus \GL(\CC^M)$-module. The super commutator is identically zero.
\end{exam}

A \emph{Lie superalgebra} $\fgtil$ is a $\ZZ/2$-graded vector space $\fgtil = \fgtil_0 \oplus \fgtil_1$ with a bilinear operation $[-, -]$ satisfying (S1) and (S2) with $[\fgtil_i, \fgtil_j] \subset \fgtil_{i+j}$. As with all superalgebras, the tensor superalgebra $\Trmtil(V)$ and the symmetric superalgebra $\Stil(V)$ are Lie superalgebras under the super commutator.

\begin{exam}
    The \textit{free Lie superalgebra} generated by $V = V_0 \oplus V_1 = \CC^N \oplus \CC^M$ is the Lie superalgebra $\Ltil(V)$ in $\Trmtil(V)$ generated by $V$. As with $\Trmtil(V)$, it is a bigraded and $\ZZ/2$-graded $\GL(\CC^N) \oplus \GL(\CC^M)$-module. Concretely, $\Ltil(V) \subset \Trmtil(V)$ is given by
    \begin{align*}
        \Ltil_{n, m}(V) &\coloneqq \text{span of iterated bracketings with $n$ terms from $V_0$, $m$ terms from $V_1$} \\
        \Ltil_0(V) &\coloneqq \text{span of iterated bracketings with evenly many terms from $V_1$} \\
        \Ltil_1(V) &\coloneqq \text{span of iterated bracketings with oddly many terms from $V_1$}.
    \end{align*}
    Note that $\Ltil_{0, 0}(V) = 0$, $\Ltil_{1, 0}(V) = V_0$, and $\Ltil_{0, 1}(V) = V_1$.
\end{exam}

Petrogradsky \cite{Pet} found the doubly graded Hilbert series of $\Ltil(V)$, generalizing what is known in the classical case as \emph{Witt's formula}, which in turn yields a formula for the graded character of $\Ltil(V)$ as a $\GL(\CC^N) \oplus \GL(\CC^M)$-module.
\begin{thm}[\cite{Pet}]\label{thm:free-super-lie-char}
    The bihomogeneous $\GL(\CC^N) \oplus \GL(\CC^M)$-character of $\Ltil(V)$ is
    \[
    \grch(\Ltil(V);\xb_N,\yb_M,q,t) = -\sum_{d=1}^\infty \frac{\mu(d)}{d} \ln(1-(q^dp_d(\xb_N) - t^dp_d(-\yb_M))).
    \]
\end{thm}
The above result of Petrogradsky allows us to determine a character formula for the doubly graded pieces $\Ltil_{n,m}(V)$ of $\Ltil(V)$.
\begin{thm}\label{thm:super-bi-brandt}
    For $(n,m) \neq (0,0)$, the $\GL(\CC^N) \oplus \GL(\CC^M)$-character of $\Ltil_{n,m}(V)$ is
    \[
    \ch(\Ltil_{n,m}(V);\xb_N,\yb_M) = \frac{1}{n+m} \sum_{d \mid \gcd(m,n)} (-1)^\frac{m}{d} \mu(d) \binom{\frac{n+m}{d}}{\frac{m}{d}} p_d(\xb_N)^\frac{n}{d}p_d(-\yb_M)^\frac{m}{d}.
    \]
\end{thm}
\begin{proof}
    We have
    \begin{align*}
    \grch(\Ltil(V);&\xb_N,\yb_M,q,t) = -\sum_{d=1}^\infty \frac{\mu(d)}{d} \ln(1-(q^dp_d(\xb_N) - t^dp_d(-\yb_M))) \\
    &= \sum_{d=1}^\infty \frac{\mu(d)}{d} \sum_{s=1}^\infty \frac{(q^dp_d(\xb_N) - t^dp_d(-\yb_M))^s}{s} \\
    &= \sum_{d=1}^\infty \frac{\mu(d)}{d} \sum_{s=1}^\infty \frac{1}{s} \sum_{k=0}^s {\binom{s}{k}}(q^dp_d(\xb_N))^k(-t^dp_d(-\yb_M))^{s-k} \\
    &= \sum_{d=1}^\infty \sum_{\substack{0 \le k \le s < \infty, \\ (k,s) \neq (0,0)}} \frac{\mu(d)}{d} \frac{1}{s} {\binom{s}{k}} (q^dp_d(\xb_N))^k(-t^dp_d(-\yb_M))^{s-k} \\
    &= \sum_{d=1}^\infty \sum_{\substack{n : d \mid n, \\ m : d \mid m, \\ (n,m) \neq (0,0)}} \frac{\mu(d)}{n+m} \binom{\frac{n+m}{d}}{\frac{m}{d}} p_d(\xb_N)^\frac{n}{d}(-p_d(-\yb_M))^\frac{m}{d} q^nt^m.
    \end{align*}
    As $\ch(\Ltil_{n,m}(V);\xb_N,\yb_M) = [q^nt^m]\grch(\Ltil(V);\xb_N,\yb_M,q,t)$, the result follows.
\end{proof}
As a corollary to \Cref{thm:super-bi-brandt}, we obtain the super generalization of Brandt's formula, \Cref{thm:super-brandt}. Recall that we set $N=M$ and use the diagonal action here.
\begin{proof}[Proof (of \Cref{thm:super-brandt}).]
By \Cref{thm:super-bi-brandt}, we have
\begin{align*}
\ch(\Ltil_{n,m}(V);\xb_N) &= \frac{1}{n+m} \sum_{d \mid \gcd(m,n)} (-1)^\frac{m}{d} \mu(d) \binom{\frac{n+m}{d}}{\frac{m}{d}} p_d(\xb_N)^\frac{n}{d}p_d(-\xb_N)^\frac{m}{d} \\
&= \frac{1}{n+m} \sum_{d \mid \gcd(n,m)} (-1)^{m + \frac{m}{d}} \mu(d) \binom{\frac{n+m}{d}}{\frac{m}{d}} p_d(\xb_N)^\frac{n+m}{d}. \qedhere
\end{align*}
\end{proof}

\subsection{Thrall's problem for free Lie superalgebras}\label{sec:super-Thrall}
We now define the super higher Lie modules (see \Cref{def:super-Lie}) which form a $\GL(\CC^N) \oplus \GL(\CC^M)$ decomposition of the tensor algebra of a super vector space described in \Cref{thm:super-Thrall}. We then state the super generalization of Thrall's problem.

The \textit{universal enveloping superalgebra} of a Lie superalgebra $\fgtil$ is intuitively the smallest superalgebra $\Util(\fgtil)$ which $\fgtil$ embeds into as a Lie superalgebra. Formally, $\Util(\fgtil)$ is a superalgebra with a map of Lie superalgebras $\iota \colon \fgtil \to \Util(\fgtil)$ satisfying the following universal property. For any superalgebra $\Atil$, composition with $\iota$ induces a bijection
  \[ \Hom_{\textrm{SuperLie}}(\fgtil, \Atil) \cong \Hom_{\textrm{SuperAlg}}(\Util(\fgtil), \Atil). \]
This leads to a natural construction of $\Util(\fgtil)$ as a quotient of $\Trmtil(\fgtil)$, which results in a filtration $\Util(\fgtil) = \bigcup_{d \geq 0} \Util_{\leq d}(\fgtil)$. The \textit{associated graded} superalgebra is $\gr \Util(\fgtil) \coloneqq \bigoplus_{d \geq 0} \Util_{\leq d}(\fgtil)/\Util_{\leq d-1}(\fgtil)$ where $\Util_{\leq -1}(\fgtil) \coloneqq 0$. Note that $\Util(\fgtil) \cong \gr \Util(\fgtil)$ as vector spaces.

\begin{exam}\label{ex:U-of-free}
    The universal enveloping superalgebra of the free Lie superalgebra $\Ltil(V)$ is $\Trmtil(V)$. Indeed, if $\Atil$ is a superalgebra, $X$ is a basis for $V_0$, and $Y$ is a basis for $V_1$, we have natural identifications
    \begin{align*}
        \Hom_{\text{SuperAlg}}(\Util(\Ltil(V)), \Atil)
          &\cong \Hom_{\text{SuperLie}}(\Ltil(V), \Atil) \\
          &\cong \Hom_{\text{SuperSet}}(X \sqcup Y, \Atil_0 \sqcup \Atil_1) \\
          &\cong \Hom_{\text{SuperVec}}(V_0 \oplus V_1, \Atil_0 \oplus \Atil_1) \\
          &\cong \Hom_{\text{SuperAlg}}(\Trmtil(V), \Atil).
    \end{align*}
    Hence $\Hom_{\text{SuperAlg}}(\Util(\Ltil(V)), -) \cong \Hom_{\text{SuperAlg}}(\Trmtil(V), -)$, so by Yoneda's lemma,
      \[ \Util(\Ltil(V)) \cong \Trmtil(V). \]
\end{exam}

The Poincar\'e--Birkhoff--Witt theorem for Lie superalgebras reads as follows. A standard consequence is that the map $\iota$ is an embedding, though we will only apply the result when $\fgtil = \Ltil(V)$, in which case injectivity is obvious by \Cref{ex:U-of-free}.

\begin{thm}[Super PBW]
   Let $\fgtil$ be a Lie superalgebra. Then there is a natural isomorphism of superalgebras
     \[ \Stil(\fgtil) \cong \gr \Util(\fgtil). \]
\end{thm}

We are nearly ready to state the super analogue of Thrall's decomposition of the tensor algebra. We first define analogues of Lie modules, which will arise naturally in the upcoming proof.

\begin{defi}\label{def:super-Lie}
    For a vector space $W$ and $j \in \ZZ_{\geq 0}$, let
    \[
    \Gamma_j(W) =
    \begin{cases}
        S(W) & \text{if $j$ is even} \\
        \bigwedge(W) & \text{if $j$ is odd}.
    \end{cases}
    \]
    Given a $\ZZ_{\ge 0}$-valued matrix $A = (a_{i,j})_{i,j \ge 0}$ with finite support and $a_{0,0} = 0$, the \emph{super Lie module} $\Ltil_A$ is:
    \begin{equation}\label{eq:super-higher-lie}
    \Ltil_A = \bigotimes_{i,j \ge 0} \Gamma_j^{a_{i,j}}(\Ltil_{i,j}).
    \end{equation}
    Note that for $j$ odd, the exterior power $\Gamma_j^a(W)$ is zero unless $a \leq \dim W$.
\end{defi}

\begin{thm}[Super Thrall decomposition]\label{thm:super-Thrall}
    Let $V = V_0 \oplus V_1 = \CC^N \oplus \CC^M$. Then
      \[ \Trmtil(V) = \bigoplus_A \Ltil_A \]
    as $\GL(\CC^N) \oplus \GL(\CC^M)$-modules, where the sum is over $\ZZ_{\geq 0}$-valued matrices $A = (a_{i,j})_{i,j \geq 0}$ with finite support and $a_{0,0} = 0$.
\end{thm}

\begin{proof}
    By \Cref{ex:U-of-free}, $\Trmtil(V) = \Util(\Ltil(V))$. Hence $\gr \Util(\Ltil(V)) = \Util(\Ltil(V))$. The super PBW theorem and routine properties now gives
    \begin{align*}
        \Trmtil(V)
          &= \Util(\Ltil(V)) = \gr \Util(\Ltil(V)) = \Stil(\Ltil(V)) = S(\Ltil(V)_0) \otimes \bigwedge(\Ltil(V)_1) \\
          &= S\left(\bigoplus_{n,m} \Ltil_{n, 2m}\right) \otimes \bigwedge\left(\bigoplus_{n,m} \Ltil_{n, 2m+1}\right)
           = \bigotimes_{n,m} S\left(\Ltil_{n, 2m}\right) \otimes \bigotimes_{n, m} \bigwedge\left(\Ltil_{n, 2m+1}\right) \\
          &= \bigotimes_{i, j} \Gamma_j\left(\Ltil_{i, j}\right)
           = \bigotimes_{i, j} \left( \bigoplus_{a_{i, j} \geq 0} \Gamma_j^{a_{i, j}}\left(\Ltil_{i, j}\right) \right)
           = \bigoplus_{A = (a_{i, j} \geq 0)} \bigotimes_{i, j} \Gamma_j^{a_{i, j}}\left(\Ltil_{i, j}\right) \\
          &= \bigoplus_{A = (a_{i, j} \geq 0)} \Ltil_A.\qedhere
    \end{align*}
\end{proof}

Now, if $N = M$, then $\Ltil$ inherits the structure of a $\GL(\CC^N)$-module under the diagonal inclusion $\GL(\CC^N) \hookrightarrow \GL(\CC^N) \times \GL(\CC^N)$. Then Thrall's problem may be generalized to the free Lie superalgebra as follows.
\begin{prob}[Super Thrall's problem]
  For $A = (a_{i, j} \geq 0)_{i, j}$ and $\lambda \in \mathrm{Par}$, what is the multiplicity of the Schur module $V^\lambda$ in the super higher Lie module $\Ltil_A$ where $N=M \to \infty$?
\end{prob}
As in the classical case, it would be sufficient to determine these multiplicities for the modules $S^a(\Ltil_{n,2m})$ and $\bigwedge^b(\Ltil_{n,2m+1})$ for general $a,b,n,m \in \ZZ_{\ge 0}$, by the Littlewood--Richardson rule. As a result, the corresponding problem for the modules $\Ltil_{n,m}$ is of particular importance.

\section{Schur--Weyl duals as induced characters}\label{sec:super-induced}
In this section we prove \Cref{thm:super-induced}, which identifies the Schur--Weyl dual of $\Ltil_{n,m}$. We first describe a certain representation of $C_{n+m}$ and then prove that, upon inducing up to $S_{n+m}$, the Frobenius character of the resulting representation agrees with $\Ch(\Ltil_{n,m})$.

Recall from above that $C_r \le S_r$ is the cyclic group of order $r$, generated by the cycle $\pi_r = (1 2 \cdots r) \in S_r$, and $\omega_r$ denotes a primitive $r$-th root of unity. The irreducible characters of $C_r$ are $\chi^1, \ldots, \chi^r$, given by $\chi^d(\pi_r) = \omega_r^d$ for $1 \le d \le r$.

The group $C_r$ acts naturally on $[r]$ by cyclic rotation. Fixing $n,m \ge 0$ with $(n,m) \neq (0,0)$, let $\binom{[n+m]}{m} = \{ S \subseteq [n+m] : |S| = m \}$. Then the action of $C_{n+m}$ on $[n+m]$ extends naturally to an action on $\binom{[n+m]}{m}$:
\[
\pi_{n+m}^k \cdot \{ i_1, \ldots, i_m \} = \{ i_1 + k\pmod{n+m}, \ldots, i_m + k\pmod{n+m} \}
\]
for any $k \ge 1$. Let $\chi^\mathrm{cyc} : C_{n+m} \to \CC$ denote the character of this latter action, so that
\[
\chi^\mathrm{cyc}(\pi_{n+m}^k) = |\{ S \in \binom{[n+m]}{m} : \pi_{n+m}^k \cdot S = S \}|.
\]

\begin{lemma}\label{cyc-char}
For any $n,m \ge 0$ with $(n,m) \neq 0$ and $1 \le k \le n+m$, let $d = \frac{n+m}{\gcd(k,n+m)}$. Then
\[
\chi^\mathrm{cyc}(\pi_{n+m}^k) = \begin{cases}
\binom{\frac{n+m}{d}}{\frac{m}{d}} & \text{if } d \mid m \\
0 & \text{if } d \nmid m.
\end{cases}
\]
\end{lemma}
\begin{proof}
Note that $\pi_{n+m}^k$ is a product of $\gcd(k,n+m) = \frac{n+m}{d}$ disjoint cycles, each of length $d$. Let $[n+m] = C_1 \sqcup \cdots \sqcup C_{\frac{n+m}{d}}$ denote the partition of $[n+m]$ given by the decomposition of $\pi_{n+m}^k$ into disjoint cycles. Now, consider the action of $\pi_{n+m}^k$ on $\binom{[n+m]}{m}$. A set $S \subseteq [n+m]$ is fixed by $\pi_{n+m}^k$ if and only if, for each $1 \le j \le \gcd(k,n+m)$, either $S \cap C_j = C_j$ or $S \cap C_j = \varnothing$. Thus the $m$-subsets of $[n+m]$ fixed by $\pi_{n+m}^k$ may be constructed by choosing $\frac{m}{d}$ of the sets $C_1,\ldots,C_{\frac{n+m}{d}}$. This is impossible if $d \nmid m$, and the result follows.
\end{proof}

For any subgroup $H \le S_r$ and $\chi$ a representation of $H$, the Frobenius characteristic of $\chi \uparrow_H^{S_r}$ may be easily expressed in the power sum basis (see e.g. \cite[Theorem~13]{Swa}). Here the cycle type of $\pi_{n+m}^k$ depends only on $\gcd(k,n+m)$, giving:
\begin{lemma}\label{induced-rep-char}
For any representation $\chi$ of $C_r$, the Frobenius characteristic of $\chi \uparrow_{C_r}^{S_r}$ is given by:
\[
\FrobCh(\chi \uparrow_{C_r}^{S_r}) = \frac{1}{r} \sum_{d \mid r} \sum_{k : \gcd(k,r) = \frac{r}{d}} \chi(\pi_r^k) p_d(\mathbf{x})^{\frac{r}{d}}.
\]
\end{lemma}
We now proceed to the proof of \Cref{thm:super-induced}.
\begin{proof}[Proof (of \Cref{thm:super-induced}).]
Throughout, let $\pi = \pi_{n+m}$ and $\omega = \omega_{n+m}$. 

First assume $m$ is odd, so that for any $d \mid m$, $m+\frac{m}{d}$ is always even. Then by \Cref{induced-rep-char}, we have:
\[
\FrobCh((\chi^\mathrm{cyc} \otimes \chi^1)\uparrow_{C_{n+m}}^{S_{n+m}}) = \frac{1}{n+m} \sum_{d \mid n+m} \sum_{\substack{k\\\gcd(k,n+m) = \frac{n+m}{d}}} \chi^\mathrm{cyc}(\pi^k) \chi^1(\pi^k) p_d(\mathbf{x})^{\frac{n+m}{d}},
\]
which by \Cref{cyc-char} simplifies to:
\[
\FrobCh((\chi^\mathrm{cyc} \otimes \chi^1)\uparrow_{C_{n+m}}^{S_{n+m}}) = \frac{1}{n+m} \sum_{\substack{d \mid n+m, \\ d \mid m}} \binom{\frac{n+m}{d}}{\frac{m}{d}} \sum_{\substack{k \\ \gcd(k,n+m) = \frac{n+m}{d}}} \chi^1(\pi^k) p_d(\mathbf{x})^{\frac{n+m}{d}}
\]
Note that $d \mid n+m$ and $d \mid m$ if and only if $d \mid \gcd(n,m)$. Furthermore, if $\gcd(k,n+m) = \frac{n+m}{d}$, then $\chi^1(\pi^k) = \omega^k$ is a primitive $d$-th root of unity, and all primitive $d$-th roots of unity are obtained in this way. The M{\"o}bius function $\mu(d)$ is given as the sum of the primitive $d$-th roots of unity. Thus
\[
\FrobCh((\chi^\mathrm{cyc} \otimes \chi^1)\uparrow_{C_{n+m}}^{S_{n+m}}) = \frac{1}{n+m} \sum_{\substack{d \mid \gcd(n,m)}} \binom{\frac{n+m}{d}}{\frac{m}{d}} \mu(d) p_d(\mathbf{x})^{\frac{n+m}{d}} = \Ch(\Ltil_{n,m}),
\]
where the last equality follows from \Cref{thm:super-brandt}.

Now assume $m$ is even. Again applying \Cref{cyc-char} and \Cref{induced-rep-char}, we have
\begin{align*}
\FrobCh((\chi^\mathrm{cyc} &\otimes \chi^{\frac{m}{2}+1})\uparrow_{C_{n+m}}^{S_{n+m}}) \\
&= \frac{1}{n+m} \sum_{d \mid \gcd(n,m)} \binom{\frac{n+m}{d}}{\frac{m}{d}} \sum_{k : \gcd(k,n+m) = \frac{n+m}{d}} \chi^{\frac{m}{2}+1}(\pi^k) p_d(\mathbf{x})^{\frac{n+m}{d}}.
\end{align*}
Observe for any $d \mid m$ that $\frac{m}{d}$ is even if and only if $d \mid \frac{m}{2}$, and recall from above that if $\gcd(k,n+m) = \frac{n+m}{d}$, then $\omega^k$ is a primitive $d$-th root of unity. Thus in the case that $\frac{m}{d}$ is even, we have $(\omega^k)^{\frac{m}{2}} = 1$ since $d \mid \frac{m}{2}$, which equals $(-1)^{m+\frac{m}{d}}$ since $m$ and $\frac{m}{d}$ are even. On the other hand, if $\frac{m}{d}$ is odd, then $d$ must necessarily be even since $m$ is, so that
\[
(\omega^k)^{\frac{m}{2}} = ((\omega^k)^{\frac{d}{2}})^{\frac{m}{d}} = (-1)^{\frac{m}{d}} = (-1)^{m + \frac{m}{d}} = -1.
\]
Hence
\[
\chi^{\frac{m}{2}+1}(\pi^k) = (\omega^k)^{\frac{m}{2}}\omega^k = (-1)^{m+\frac{m}{d}}\omega^k.
\]
It then follows that
\begin{align*}
\FrobCh((\chi^\mathrm{cyc} &\otimes \chi^{\frac{m}{2}+1})\uparrow_{C_{n+m}}^{S_{n+m}}) \\
&= \frac{1}{n+m} \sum_{d \mid \gcd(n,m)} (-1)^{m + \frac{m}{d}} \mu(d) \binom{\frac{n+m}{d}}{\frac{m}{d}} p_d(\xb)^\frac{n+m}{d} = \Ch(\Ltil_{n,m}). \qedhere
\end{align*}
\end{proof}

\section{Super tableau combinatorics}\label{sec:super-tab}
In this section we prove \Cref{thm:maj-neg-hook}, which is the key ingredient in our proof of \Cref{thm:super-KW}. We first recall some preliminaries on supersymmetric functions in \Cref{sec:supersym-fxns}. In \Cref{sec:super-des}, we define a new major index statistic on super tableaux and then prove \Cref{prop:q-ps} and \Cref{thm:s-ps}, which exhibit how the principal specializations of the super quasisymmetric and super Schur functions admit elegant formulae in terms of this new statistic. \Cref{thm:maj-neg-hook} then follows as an easy consequence.
\subsection{Supersymmetric functions}\label{sec:supersym-fxns}
We begin by recalling some facts about supersymmetric and super quasisymmetric functions, following the presentation in \cite{HHLRU}. Let $\yb = (y_1,y_2,\ldots)$. A power series $f(\xb;\yb) \in \Lambda(\xb) \otimes_\ZZ \Lambda(\yb)$ is \emph{supersymmetric} if performing the substitution $x_1 = t, y_1 = -t$ into $f$ results in an expression which is independent of $t$. Schur functions and power sums both admit supersymmetric analogs. The \emph{super power sum symmetric function} $\ptil_\lambda(\xb;\yb)$ indexed by $\lambda$ is:
\[
\ptil_\lambda(\xb;\yb) = \prod_{j=1}^{\ell(\lambda)} (p_{\lambda_j}(\xb) + (-1)^{\lambda_j+1}p_{\lambda_j}(\yb)).
\]

The super Schur functions were originally introduced by Berele--Regev \cite{BReg} in their study of the general linear Lie superalgebra, and they may be defined as follows. Let $\Acal_+ = \{ 1, 2, \ldots \}$ and $\Acal_- = \{ \overline{1}, \overline{2}, \ldots \}$. Endow the alphabet $\Acal = \Acal_+ \sqcup \Acal_-$ with the following total order:
\[
\Acal = \{ 1 < \overline{1} < 2 < \overline{2} < \cdots \}.
\]
We define a projection $\Acal \to \Acal_+$ by $i \mapsto i$ and $\overline{i} \mapsto i$ for all $i \ge 1$. Note that if $a_1 \le \cdots \le a_n$ is a weakly increasing sequence in $\Acal$ and $a_j \mapsto b_j$ under this projection, then $b_j \le a_j$ for all $j$ and $b_1 \le \cdots \le b_n$.
 
\begin{defi}[\cite{HHLRU}, Eq.~(23)]
For $n \ge 2$ and $D \subseteq [n-1]$, the \emph{super quasisymmetric function} $\Qtil_{n,D}(\xb;\yb)$ is given by
\[
\Qtil_{n,D}(\xb;\yb) = \sum_{\substack{a_1 \le a_2 \le \cdots \le a_n, \\ a_i = a_{i+1} \in \Acal_+ \Rightarrow i \not\in D, \\ a_i = a_{i+1} \in \Acal_- \Rightarrow i \in D}} z_{a_1} z_{a_2} \cdots z_{a_n},
\]
where $a_1 \le \cdots \le a_n$ is a weakly increasing sequence in $\Acal$, and $z_a = x_a$ for $a \in \Acal_+, z_b = y_b$ for $b \in \Acal_-$.
\end{defi}
As in the classical case, the super Schur function $\stil_\lambda(\xb;\yb)$ is given as a sum of super quasisymmetric functions:
\begin{defi}[\cite{HHLRU}, Prop.~2.4.2]\label{prop:super-s-to-q}
    The \emph{super Schur function} $\stil_\lambda(\xb;\yb)$ is given in terms of super quasisymmetric functions by
    \[
    \stil_\lambda(\xb;\yb) = \sum_{T \in \SYT(\lambda)} \Qtil_{|\lambda|,\Des(T)}(\xb;\yb).
    \]
\end{defi}

Schur functions and power sum symmetric functions are related via the \emph{Cauchy identity}, which extends to the supersymmetric setting as follows.
\begin{thm}[\cite{BRem}, Cor.~10(a)]\label{lem:s-p}
    For all $n \ge 1$,
    \[
    \sum_{\lambda \vdash n} s_\lambda(\xb) \stil_\lambda(\xb;\yb) = \sum_{\lambda \vdash n} \frac{1}{z_\lambda} p_\lambda(\xb) \ptil_\lambda(\xb;\yb),
    \]
    where if $\lambda = (1^{a_1}2^{a_2} \cdots)$, then $z_\lambda \coloneqq 1^{a_1}2^{a_2} \cdots a_1!a_2! \cdots$.
\end{thm}

\subsection{A major index statistic on super tableaux}\label{sec:super-des} We now define a new major index statistic in \Cref{def:super-maj} on objects called \emph{standard super tableaux}, and show how particular specializations of super quasisymmetric and super Schur functions may be written in terms of this statistic in \Cref{prop:q-ps} and \Cref{thm:s-ps}.
\begin{defi}
    A \emph{standard super tableau} of shape $\lambda \vdash n$ is a map $\mathcal{T} : \lambda \to \Acal$ that is strictly increasing along the rows and columns of $\lambda$, and contains exactly one of $i$ or $\overline{i}$ for each $i = 1, 2, \ldots, n$. Let $\SYT_\pm(\lambda)$ denote the set of standard super tableaux of shape $\lambda$, so that $|\SYT_\pm(\lambda)| = 2^n|\SYT(\lambda)|$. For $\mathcal{T} \in \SYT_\pm(\lambda)$, we let $\Neg(\mathcal{T}) \coloneqq \{ i \in \Acal_+ : \overline{i} \in \mathcal{T} \}$ and $\negg(\mathcal{T}) \coloneqq |\Neg(\mathcal{T})|$.
\end{defi}
